\section*{Methodology}

I would estimate a difference-in-differences model by OLS to compare investment changes in exposed and unexposed firms. Here, $i$ identifies firms and $t$ identifies quarters. The dependent variable is investment $I_{it}$, defined above. The explanatory variables combine exposure $E_i$, initial cash holdings $C_i$ and a dummy $\mathrm{Post}_t$ equal to one from 2024Q4 onwards and zero before Helene:

\begin{equation*}
\begin{aligned}
I_{it} ={}& \alpha_i + \lambda_t + \beta_1(E_i \times \mathrm{Post}_t) \\
&+ \beta_2(\mathrm{Post}_t \times C_i)
 + \beta_3(E_i \times \mathrm{Post}_t \times C_i) + \varepsilon_{it}.
\end{aligned}
\end{equation*}

Firm fixed effects ($\alpha_i$) account for persistent differences between firms. Quarter fixed effects ($\lambda_t$) capture shocks shared across firms in each quarter, and $\varepsilon_{it}$ is the error term. The simpler terms $E_i$, $C_i$ and $E_i \times C_i$ are absorbed by firm effects because they do not change over time. Quarter effects absorb $\mathrm{Post}_t$. Thus, all lower-order terms underlying the three-way interaction are included or absorbed.

The term $\mathrm{Post}_t \times C_i$ allows cash to be associated with investment changes even among unexposed firms. Consequently, the exposure-related change at cash level $C_i$ is $\beta_1 + \beta_3 C_i$. The coefficient $\beta_3$ answers the research question: a ten-percentage-point difference in cash/assets implies a difference of $0.100\beta_3$ percentage points in the quarterly investment response.

I would test $H_0: \beta_3 \leq 0$ against $H_1: \beta_3 > 0$ using a one-sided t-test with standard errors clustered by firm, and report a 95\% confidence interval. The comparison assumes that, without Helene, exposed and unexposed firms would have followed comparable investment trends at similar cash levels. Pre-storm plots for both groups, separated by initial cash levels, would help assess this assumption without proving it.
